% answers_short.tex —— 2010 年·普通高等学校招生考试·四川卷·理科 · 简短答案速查
% 由 tools/make_short_answers.py 自 answers.tex 的 \gksolhead/\gkanswer 提取，
% 并与 2010/2010_数学真题_四川理_答案要点.md 逐题比对后生成，勿手改（改源 MD 后重跑）。

\begin{gknotice}
\textbf{使用说明}：本节只给最终答案便于快速核对，完整推导见解析版。答案为编写者独立推导，\textbf{非官方标准答案}。
\end{gknotice}

\gkpart{选择题}
\begin{center}\begin{tabular}{ccccccccccccc}
\toprule
题号 & 1 & 2 & 3 & 4 & 5 & 6 & 7 & 8 & 9 & 10 & 11 & 12 \\
\midrule
答案 & A & D & C & A & C & C & B & B & D & C & A & B \\
\bottomrule
\end{tabular}\end{center}

\gkpart{填空题}
\begin{itemize}[leftmargin=2.6em]
  \item[\textbf{13.}] $-\dfrac{160}{x}$。
  \item[\textbf{14.}] $2\sqrt3$。
  \item[\textbf{15.}] $\dfrac{\sqrt3}{4}$。
  \item[\textbf{16.}] ①②。
\end{itemize}

\gkpart{解答题}
\begin{itemize}[leftmargin=2.6em]
  \item[\textbf{17.}] （1）$\dfrac{25}{216}$；（2）分布列见解析，$E\xi=\dfrac12$。
  \item[\textbf{18.}] （1）证明；（2）$\arccos\dfrac13$；（3）$\dfrac{1}{24}$。
  \item[\textbf{19.}] （1）证明；（2）$-\dfrac{\sqrt{10}}{10}$。
  \item[\textbf{20.}] （1）$x^2-\dfrac{y^2}{3}=1$；（2）以 $MN$ 为直径的圆过点 $F$。
  \item[\textbf{21.}] （1）$a_3=6$，$a_5=20$；（2）$b_n=8n-2$，公差 $8$；（3）见解析。
  \item[\textbf{22.}] （1）$5\leqslant t\leqslant32$；（2）证明；（3）$\left|\sum\limits_{k=1}^nf(k)-n\right|<4$。
\end{itemize}
